% Figure : chaîne de calcul de la matrice de couplage électron-défaut M jusqu'à sa troncature (fig:chaine_M)
% Même style que les diagrammes SCF (scf.tex, ks_pw.tex, dfpt.tex) : boîtes grises, numéros en haut à gauche,
% boîtes latérales pour les quantités préalables, cadres pointillés pour regrouper les étapes.
% Deux entrées (super-cellules, cellule unitaire), deux branches (locale, non locale), puis la chaîne en base de Wannier.
% La sortie (M tronquée) est l'entrée de la figure analogue de la section 4.2.
\begin{tikzpicture}[
    box/.style={rectangle, draw=black, rounded corners, fill=gray!15,
                text width=6.0cm, inner xsep=0.5cm, inner ysep=5pt,
                minimum height=1.1cm, align=center, font=\linespread{1}\selectfont\small},
    inbox/.style={box, text width=4.8cm},
    sidebox/.style={rectangle, draw=black, rounded corners, fill=gray!15,
                text width=2.7cm, inner ysep=4pt, minimum height=0.85cm,
                align=center, font=\linespread{1}\selectfont\small},
    arrow/.style={->, >=stealth, thick},
    num/.style={font=\scriptsize\bfseries, anchor=north west, inner sep=2pt},
    label/.style={font=\scriptsize},
    cadre/.style={dashed, gray, rounded corners},
    titre/.style={font=\small\itshape, gray},
]

% ---- Entrées ----
\node[inbox] (sc) at (-3.2, 0)
    {\texttt{QuantumESPRESSO} \\ Super-cellules $N\times N$ parfaite et défectueuse \\[2pt]
     $V^\text{L}, \hat V^\text{NL}$, $\bm{\tau}_\alpha$};
\node[num] at (sc.north west) {(1)};

\node[inbox] (uc) at (3.2, 0)
    {\texttt{QuantumESPRESSO} \\ Cellule unitaire grille $pN\times pN$ \\[2pt]
     $C_{n\kv}(\Gv)$, $\varepsilon_{n\kv}$};
\node[num] at (uc.north west) {(2)};

% ---- Base de Bloch : parties locale et non locale ----
\node[inbox] (pad) at (-3.2, -2.3)
    {Remplissage par zéros\\($pN\times pN$) puis FFT\\[2pt]
     $\widetilde{V}^\text{L}_\text{ed}(\qv+\Gv)$};
\node[num] at (pad.north west) {(3)};

\node[inbox] (ML) at (-3.2, -4.4)
    {Partie locale\\éq.~\eqref{eqn:ML_final}\\[2pt]
     $M^\text{L}_{mn}(\kv',\kv)$};
\node[num] at (ML.north west) {(4)};

\node[inbox] (MNL) at (3.2, -4.4)
    {Partie non locale\\éq.~\eqref{eqn:MNL_final}\\[2pt]
     $M^\text{NL}_{mn}(\kv',\kv)$};
\node[num] at (MNL.north west) {(5)};

\node[box] (M) at (0, -6.4)
    {$M_{mn}=M^\text{L}_{mn}+M^\text{NL}_{mn}$\\[2pt]
     {\scriptsize grille  $pN\times pN$}};
\node[num] at (M.north west) {(6)};

% ---- Base de Wannier ----
\node[box] (MW) at (0, -8.8)
    {Base de Wannier, éq.~\eqref{eqn:M_wannier}\\[2pt]
     $M_{ij}(\Rv',\Rv)$ sur les $(pN)^2$ cellules};
\node[num] at (MW.north west) {(7)};

\node[box] (align) at (0, -10.6)
    {Alignement, éq.~\eqref{eqn:alignement_W}\\[2pt]
     $-\Delta V^{(N)}_\text{PA}\,\delta_{ij}\delta_{\Rv\Rv'}$ sur les $N^2$ cellules seulement};
\node[num] at (align.north west) {(8)};

\node[box] (rec) at (0, -12.2)
    {Recentrage sur la lacune\\éq.~\eqref{eqn:M_recentrage}};
\node[num] at (rec.north west) {(9)};

\node[box] (trunc) at (0, -13.8)
    {Troncature : $R, R'\le R_\text{cut}$\\[2pt]
     $M_{ij}(\Rv',\Rv)$ tronquée};
\node[num] at (trunc.north west) {(10)};

% ---- Quantités préalables (boîtes latérales) ----
\node[sidebox] (wan) at (5.6, -8.8)
    {\texttt{Wannier90}\\$U^{(\kv)}$, $H(\Rv)$};

\node[sidebox] (PA) at (-5.6, -10.6)
    {$\Delta V^{(N)}_\text{PA}$, éq.~\eqref{eqn:VPA}\\{\scriptsize loin de la lacune}};

% ---- Flèches ----
\draw[arrow] (sc) -- (pad);
\draw[arrow] (pad) -- (ML);
\draw[arrow] (uc) -- (MNL);
\draw[arrow] (sc.south east) -- ([xshift=-1.6cm]MNL.north);
\draw[arrow] (uc.south west) -- ([xshift=1.6cm]ML.north);
\draw[arrow] (ML.south) -- ([xshift=-1.8cm]M.north);
\draw[arrow] (MNL.south) -- ([xshift=1.8cm]M.north);
\draw[arrow] (M) -- (MW);
\draw[arrow] (MW) -- (align);
\draw[arrow] (align) -- (rec);
\draw[arrow] (rec) -- (trunc);

\draw[arrow] (uc.east) -- ++(1.3, 0) |- (wan.east);
\draw[arrow] (wan) -- (MW);
\draw[arrow] (sc.west) -- ++(-1.3, 0) |- (PA.west);
\draw[arrow] (PA) -- (align);

% ---- Sortie ----
\draw[arrow] (trunc.south) -- ++(0, -1.0)
    node[below, font=\linespread{1}\selectfont\small, align=center]
    {Interpolation et \\ quantités physiques\\
      (section~\ref{sec:mbt_défaut})};

% ---- Cadres : base de Bloch et base de Wannier ----
\draw[cadre] (-6.35, -1.35) rectangle (6.35, -7.25);
\node[titre, anchor=north east] at (6.3, -1.4) {Base de Bloch};

\draw[cadre] (-3.85, -7.55) rectangle (3.85, -14.6);
\node[titre, anchor=north east] at (3.8, -7.6) {Base de Wannier};

\end{tikzpicture}
